% BHU PYQ -- Manually transcribed & error-corrected
\documentclass[11pt,a4paper]{article}

\usepackage[a4paper,top=2.5cm,bottom=2.5cm,left=2.8cm,right=2.8cm,headheight=14pt]{geometry}
\usepackage{amsmath,amssymb}
\usepackage[T1]{fontenc}
\usepackage[utf8]{inputenc}
\usepackage{lmodern}
\usepackage{microtype}
\usepackage{fancyhdr}
\usepackage{enumitem}
\usepackage{hyperref}
\usepackage{lastpage}
\usepackage{parskip}

\hypersetup{colorlinks=true,urlcolor=black,linkcolor=black,
  pdftitle={CHB-608: Atomic & Molecular Structure -- BHU PYQ},pdfauthor={Banaras Hindu University}}

\pagestyle{fancy}\fancyhf{}
\renewcommand{\headrulewidth}{0.4pt}\renewcommand{\footrulewidth}{0.4pt}
\fancyhead[L]{\small   \textit{Banaras Hindu University}}
\fancyhead[R]{\small   \textit{CHB-608: Atomic \& Molecular Structure}}
\fancyfoot[C]{\small Page  \thepage\ of \pageref{LastPage}}


\newlist{parts}{enumerate}{2}
\setlist[parts,1]{label=(\alph*),leftmargin=*,itemsep=5pt,topsep=3pt}
\setlist[parts,2]{label=(\roman*),leftmargin=*,itemsep=3pt,topsep=2pt}

\newcommand{\pts}[1]{\hfill   \textit{[#1]}}


\begin{document}


\begin{center}
  {\large\bfseries BANARAS HINDU UNIVERSITY}\\[2pt]
  {
\normalsize B.Sc. (Hons.) Semester VI Examination, 2013-14}\\[6pt]
  \rule{0.8\linewidth}{0.5pt}\\[6pt]
  {\large\bfseries Paper No. CHB-608: Atomic and Molecular Structure and Spectroscopic Techniques}\\[4pt]
  {
\normalsize Time: 3 Hours\hfill Maximum Marks: 70}
\end{center}
\rule{\linewidth}{0.4pt}
\small



\noindent   \textit{Instructions: Answer five questions in total, including Question~1 which is compulsory.}

\normalsize
\bigskip




\noindent   \textbf{Question 1.} Answer all parts of the following: \pts{5 $ \times$ 2 = 10}

\begin{parts}
  \item Sketch the symmetry elements present in: (i) $ \text{H}_2 \text{O}$ and (ii) $ \text{NH}_3$.
  \item Write the radial, angular wave equations for the Hydrogen atom.
  \item What are the basic principles of the Valence Bond Theory?
  \item What is Fermi resonance in IR spectra?
  \item What are metastable ions in mass spectrometry?
\end{parts}

\medskip\rule{\linewidth}{0.2pt}




\noindent   \textbf{Question 2.}

\begin{parts}
  \item Explain the molecular orbital treatment of the $ \text{H}_2$ molecule. Deduce the expressions for bonding and antibonding molecular orbitals. \pts{6}
  \item Derive Schr\"odinger's wave equation for a particle in a $3$-D box and calculate the degeneracy of energy levels. \pts{8}
\end{parts}

\medskip\rule{\linewidth}{0.2pt}




\noindent   \textbf{Question 3.}

\begin{parts}
  \item How will you distinguish the following pairs with the help of IR spectroscopy? \pts{6}
  
\begin{parts}
    \item Intermolecular vs. intramolecular hydrogen bonding (dilution effect)
    \item $ \text{o}$-Hydroxybenzoic acid and $ \text{p}$-hydroxybenzoic acid
  \end{parts}
  \item Discuss the basic principles of the origin of the $^1 \text{H}$-NMR spectrum. Why is TMS used as a reference? \pts{8}
\end{parts}

\medskip\rule{\linewidth}{0.2pt}




\noindent   \textbf{Question 4.}

\begin{parts}
  \item Define: chromophore, auxochrome, bathochromic shift, and hypsochromic shift. Give examples of each. \pts{6}
  \item Predict the masses of the ions produced in the mass spectrum of $2$-hexene by $\beta$-fragmentation. \pts{4}
  \item What are the masses of the ions produced in the mass spectra of $4$-$n$-butyltoluene by McLafferty rearrangement and benzylic fission? \pts{4}
\end{parts}

\vfill
\rule{\linewidth}{0.4pt}

\begin{center}\small
  \href{https://bhu-exam.vercel.app/}{Click here to visit our website}\\[4pt]
  \href{https://me-aryan.vercel.app/}{Click here to visit our contributor}\\[4pt]
  \href{https://quashan-me.vercel.app/}{Click here to visit our contributor}
\end{center}

\end{document}
